\documentclass[11pt]{article}

\usepackage[margin=1in]{geometry}
\usepackage[T1]{fontenc}
\usepackage[utf8]{inputenc}
\usepackage{booktabs}
\usepackage{longtable}
\usepackage{array}
\usepackage{xcolor}
\usepackage{hyperref}
\usepackage{listings}
\usepackage{amsmath}
\usepackage{siunitx}

\hypersetup{
  hidelinks,
  pdftitle={GF/PED Manual},
  pdfauthor={ORACLE Development Team}
}

\definecolor{ORACLEBlue}{HTML}{245C73}
\definecolor{ORACLESlate}{HTML}{4A5561}
\definecolor{ORACLEGold}{HTML}{C88A2D}
\definecolor{ORACLEBg}{HTML}{F7F8F6}
\emergencystretch=2em

\lstdefinestyle{oracle}{
  basicstyle=\ttfamily\small,
  breaklines=true,
  columns=fullflexible,
  keepspaces=true,
  frame=single,
  rulecolor=\color{ORACLESlate!45},
  backgroundcolor=\color{ORACLEBg},
  xleftmargin=0.5em,
  xrightmargin=0.5em
}

\newcommand{\program}{\texttt{GF/PED}}
\newcommand{\oracle}{\texttt{ORACLE}}
\newcommand{\gicforge}{\texttt{GICForge}}
\newcommand{\bohr}{\text{bohr}}

\title{\textbf{GF/PED Manual}\\
\large Wilson GF analysis, internal-coordinate force constants and potential-energy distributions in ORACLE}
\author{ORACLE Development Notes}
\date{\today}

\begin{document}
\maketitle

\begin{abstract}
\program{} is the \oracle{} workflow that transforms a Cartesian harmonic
Hessian into an internal-coordinate Wilson GF problem, solves the harmonic
vibrational eigenproblem and reports frequencies, internal force constants,
normal modes and potential-energy distributions (PEDs).  The production path
uses a frozen non-redundant \gicforge{} definition, so vibrational analysis,
Gaussian input generation and structural refinement can share the same
coordinate model.  This manual describes the theory, input contracts, command
line and GUI workflows, Pulay-style scaling, output files and validation rules.
\end{abstract}

\tableofcontents
\newpage

\section{Purpose}

The GF/PED module answers a specific question: given a Cartesian Hessian and a
non-redundant internal-coordinate basis, what are the harmonic frequencies,
normal modes, internal force constants and coordinate contributions in that
basis?  It is deliberately separate from VPT2/VCI.  VPT2/VCI works in
normal-coordinate quartic force fields; GF/PED works with a Cartesian Hessian,
a Wilson B matrix and a non-redundant coordinate model.

The preferred production workflow is \texttt{gic-gf}.  It reads a frozen
\gicforge{} schema and a Gaussian FCHK Hessian adapter, evaluates the same GICs
on the Hessian geometry or on an optional external geometry, transforms the
Cartesian Hessian to the GIC basis and solves Wilson GF.  A simpler diagnostic
workflow, \texttt{gf}, reads an FCHK file and generates a ORACLE GIC basis on
the fly; it is useful for quick checks but does not provide the same
cross-module coordinate identity as a frozen \gicforge{} schema.

\section{Theoretical Background}

Let \(x\) be the \(3N\)-dimensional Cartesian coordinate vector, \(F_x\) the
Cartesian harmonic Hessian in \(E_h/\bohr^2\), and \(m_i\) the atomic masses in
amu.  For a non-redundant internal coordinate vector \(q\), the linearized
relation between Cartesian and internal displacements is
\[
  dq = B\,dx ,
\]
where \(B=\partial q/\partial x\) is the Wilson B matrix.  The kinetic-energy
matrix in internal coordinates is
\[
  G = B\,M^{-1}\,B^T ,
\]
where \(M^{-1}\) is the diagonal Cartesian inverse-mass matrix with each atomic
mass repeated for \(x,y,z\).

The internal-coordinate Cartesian backtransform used by the code is
\[
  A = M^{-1} B^T G^+ ,
\]
where \(G^+\) is the numerical inverse or pseudoinverse of \(G\).  The internal
force-constant matrix is then
\[
  F_q = A^T F_x A .
\]
The matrix is explicitly symmetrized before the GF solution.  If Pulay-style
diagonal scaling factors \(s_i\) are supplied, the scaled internal Hessian is
\[
  (F_q^\mathrm{scaled})_{ij}
  =
  (F_q)_{ij}\sqrt{s_i s_j}.
\]
This preserves symmetry and applies off-diagonal scaling as the geometric mean
of the two diagonal factors.

The Wilson GF eigenproblem is solved in symmetric form.  Since \(G\) is
positive definite for a valid non-redundant basis, the code diagonalizes
\[
  G^{1/2} F_q G^{1/2}.
\]
The eigenvalues are converted to signed harmonic frequencies in cm\(^{-1}\).
Negative eigenvalues are retained as negative frequencies and therefore expose
transition-state or instability directions rather than being silently clipped.

The internal normal modes are back-expressed in the GIC basis, and the
potential-energy distribution for coordinate \(i\) in mode \(k\) is evaluated
from the mode vector \(l_k\) as
\[
  \mathrm{PED}_{ik}
  =
  100\,
  \frac{|l_{ik}(F_q l_k)_i|}
       {\sum_j |l_{jk}(F_q l_k)_j|}.
\]
The absolute-value normalization makes the PED table robust for coupled modes
with positive and negative cross terms, while retaining the relative
coordinate contribution pattern.

\subsection{Large-amplitude GIC subspaces}

GF/PED does not identify torsions by projecting Cartesian normal modes.  The
large-amplitude coordinates are already part of the frozen non-redundant GIC
basis supplied by NEO.  GF/PED therefore treats torsions, ring puckering,
butterfly, out-of-plane, linear-bend and special fragment/centre coordinates
as ordinary rows and columns of the same internal-coordinate \(F\) and \(G\)
matrices.

For diagnostics and later anharmonic workflows, GF/PED extracts the
corresponding GIC subsets and solves mono- or poly-dimensional Wilson GF
problems on their direct \(F_{SS}\) and \(G_{SS}\) sub-blocks.  No new
coordinates are generated and no Cartesian projection is performed.  The
report also gives the largest \(F\) and \(G\) couplings between each selected
subspace and the remaining coordinates, plus the total large-amplitude PED
contribution to every GF mode.  These diagnostics decide whether a
one-dimensional, coupled multidimensional, DVR or perturbative treatment is
appropriate downstream.
The operational isolation test uses both matrices: a large-amplitude subspace
is independent only when the relative off-subspace coupling is small in \(F\)
and in \(G\).  MATRIX therefore reports the two values separately and also
exports \(\max(F_{\mathrm{rel}},G_{\mathrm{rel}})\) as the combined
\texttt{FG} diagnostic used by downstream rotor/DVR planning.
The kinetic coefficient needed by the DVR branch is not recomputed from a
separate rotor geometry algorithm.  Since the GF calculation already owns the
Wilson \(G\) matrix in the final non-redundant GIC basis, MATRIX inverts that
global matrix once and keeps \(G^{-1}\).  The DVR plan records
\((G^{-1})_{ii}\) for each large-amplitude coordinate.  For angular coordinates
this is the effective reduced moment in the current GIC metric.  For coupled
large-amplitude treatments MATRIX stores the full selected block
\(G^{-1}_{SS}\), not only its diagonal, because the multidimensional DVR
kinetic operator needs the off-diagonal metric couplings as well.
Every torsional, ring-puckering, butterfly, out-of-plane, linear-bend and
special coordinate is retained as a candidate.  The local one-coordinate GF
frequency marks whether the candidate is active; by default the CLI uses
\texttt{250 cm-1}, configurable with
\texttt{--large-amplitude-frequency-cutoff-cm}.  Coordinates above the cutoff
are written as excluded candidates rather than removed from the diagnostics.

For ordinary torsions, GF/PED also prepares a DVR plan from the frozen topology
and synthon information.  The central bond is taken from the torsional
primitive, high-order central bonds are excluded from hindered-rotor treatment,
and the periodicity is inferred from the equivalence classes of the
substituents around the central bond.  If the coordinate is isolated according
to the combined \(F/G\) coupling diagnostic, the local curvature gives the
one-term Fourier estimate
\[
  V(\phi) = A\,[1-\cos(n(\phi-\phi_0))], \qquad
  A = F_{\phi\phi}/n^2,
\]
with \(n\) the inferred periodicity and \(\phi_0\) the current minimum.  If the
combined coupling is too large, MATRIX writes a coupled DVR recommendation
instead of a one-dimensional barrier.

\section{Input Contracts}

\subsection{Cartesian Hessian}

The canonical internal object is \texttt{oracle\_gf.HessianInput}.  It stores
atomic numbers, Cartesian coordinates in bohr, atomic masses in amu, a full
symmetric Cartesian Hessian in \(E_h/\bohr^2\), optional harmonic frequencies
from the source adapter and provenance metadata.  The preferred full file
adapter is Gaussian FCHK/FCH.  Gaussian log files can also be promoted when the
Cartesian force-constant matrix is printed in the output.

The FCHK file must contain:

\begin{itemize}
  \item atomic numbers;
  \item current Cartesian coordinates;
  \item atomic masses;
  \item Cartesian force constants;
  \item harmonic frequencies when available.
\end{itemize}

\noindent
The Gaussian log adapter reads only normalized QM output data: the archive
geometry, printed Cartesian Hessian, harmonic frequencies and printed normal
coordinates when available.  It writes the shared \texttt{\#CARTESIAN\_HESSIAN}
and \texttt{\#NORMAL\_MODES} sections.  GF/PED itself does not parse Gaussian
text, does not read a Z-matrix and does not import a ReadAllGIC block from a
\texttt{.gjf} file.

\subsection{Frozen GIC Definition}

The production branch reads a \texttt{oracle.gic.definition.v1} JSON schema
created by \texttt{gic-define}.  The schema contains primitive coordinates, the
GIC coefficient matrix, labels, names, irreducible representations, point group,
symmetrization flag and the Gaussian-readable GIC block.  GF/PED evaluates that
frozen definition; it does not rerun topology perception, pruning or symmetry
assignment.

The atom count in the schema and in the FCHK input must agree.  If an optional
geometry is supplied for the B matrix, it must describe the same atom ordering
and the same topology.  Coordinates supplied through \texttt{--geometry} are in
angstrom and are converted internally to bohr before B-matrix evaluation.

\section{Command-Line Workflows}

\subsection{Quick FCHK workflow}

The quick workflow builds a ORACLE GIC basis directly from the FCHK geometry:

\begin{lstlisting}[style=oracle]
python -m matrix gf \
  --fchk working/gauin.fchk \
  --out working/gf_ped_report.txt \
  --csv-dir working/gf_csv
\end{lstlisting}

\noindent
This branch is useful for diagnostics.  Because the GICs are generated inside
the GF run, it should not be used when the analysis must be identical to a
previous Gaussian input, SEfit/MORPHEUS refinement or GICForge manifest.

\subsection{Frozen-GIC production workflow}

The production workflow first freezes the coordinate model and then runs GF/PED:

\begin{lstlisting}[style=oracle]
python -m matrix gic-define \
  --geometry parent.xyz \
  --out working/gic_definition.json \
  --gaussian-out working/gauin

python -m matrix gic-gf \
  --schema working/gic_definition.json \
  --fchk working/gauin.fchk \
  --out working/gic_gf_ped_report.txt \
  --csv-dir working/gic_gf_csv
\end{lstlisting}

\noindent
If the B matrix must be evaluated on a geometry different from the FCHK
geometry, pass:

\begin{lstlisting}[style=oracle]
python -m matrix gic-gf \
  --schema working/gic_definition.json \
  --fchk working/gauin.fchk \
  --geometry working/current.xyz \
  --out working/gic_gf_ped_report.txt
\end{lstlisting}

\noindent
This is useful for controlled diagnostics, but production vibrational analysis
should normally evaluate the B matrix on the same equilibrium geometry used for
the Hessian.

\subsection{Gaussian ReadAllGIC log workflow}

For a Gaussian job run from MATRIX-generated symmetrized GICs, keep the
coordinate definition and the QM result separate, but start the post-run
analysis from the optimized Gaussian log.  This gives MATRIX the final
geometry, not the initial geometry from the \texttt{.gjf} input.  The GICs are
rebuilt and stored by NEO/MATRIX from that final log geometry; the same
Gaussian log then contributes Hessian and normal-mode data:

\begin{lstlisting}[style=oracle]
matrix link preprocess pyrrole.log pyrrole.xyzin --source-kind gaussian
matrix validate pyrrole.xyzin
matrix neo build pyrrole.xyzin --symmetrize
matrix gaussian promote-log-hessian pyrrole.log pyrrole.xyzin
matrix gf --xyzin pyrrole.xyzin --symmetry-blocks
\end{lstlisting}

\noindent
The \texttt{pyrrole.gjf} ReadAllGIC block is therefore not treated as an input
GIC definition, and its initial Cartesian/Z-matrix geometry is not used for the
post-optimization GF/PED analysis.  This prevents GF/PED from depending on
coordinates that might have been edited, reordered or generated by a different
Gaussian-side tool.  The only accepted GIC source is the MATRIX \texttt{\#GIC}
section produced by NEO from the final log geometry.

\noindent
Gaussian log Hessians are stored after Gaussian restores the axes to the
archive/original frame.  The MATRIX log adapter uses that archive geometry for
\texttt{\#CARTESIAN\_HESSIAN}, not the last \texttt{Standard orientation}
table.  When Gaussian prints normal coordinates, \texttt{promote-log-hessian}
also writes \texttt{\#NORMAL\_MODES} and rotates those printed mode vectors to
the archive/original frame when the geometry fit is unambiguous.  Add
\texttt{--no-normal-modes} if only the Hessian section should be refreshed.

\noindent
The GF/PED report includes two diagnostics for this workflow:

\begin{itemize}
  \item a geometry check reporting raw and optimally aligned RMS/max
  differences between the frozen \texttt{\#GIC} reference geometry and the
  Hessian geometry;
  \item a frequency check reporting mode-by-mode, RMS and maximum differences
  between MATRIX GF frequencies and the harmonic frequencies stored by the
  source adapter.
\end{itemize}

\noindent
If \texttt{--symmetry-blocks} is requested, GF/PED first checks whether the
internal \(F\) and \(G\) matrices are block diagonal in the stored irrep labels.
When off-block couplings are larger than numerical roundoff, GF/PED stops with
an error: symmetry blocks from correctly symmetrized GICs must not couple
different irreps.  Such a failure is therefore a NEO/projector regression to
fix upstream, not a GF option to bypass.  PED columns are normalized to 100\%
for every solved mode.

\section{Pulay-Style Scaling}

GF/PED accepts optional diagonal scaling factors in the GIC basis.  Scaling is
applied after transformation of the Cartesian Hessian to internal coordinates
and before the Wilson GF eigenproblem.  A scaling file is line-oriented text or
CSV:

\begin{lstlisting}[style=oracle]
# selector factor
default 1.000
GIC003 0.980
A1Str0001,0.995
5=1.020
class CH_stretches 0.970 R(1,6)|R(2,7)|R(3,8)
\end{lstlisting}

\noindent
Selectors may be:

\begin{itemize}
  \item \texttt{default}, \texttt{all} or \texttt{*};
  \item a one-based index such as \texttt{5};
  \item a \texttt{GICnnn} label;
  \item an exact GIC name;
  \item an exact label;
  \item an unambiguous substring of a name or label.
\end{itemize}

\noindent
Inline records are supplied with repeated \texttt{--scale} options:

\begin{lstlisting}[style=oracle]
python -m matrix gf \
  --xyzin working/molecule.xyzin \
  --fchk working/gauin.fchk \
  --scale-file working/pulay_scale.txt \
  --scale "GIC003=0.980" \
  --scale "default=1.000" \
  --scale-class "CH_stretches:0.970:R(1,6)|R(2,7)|R(3,8)"
\end{lstlisting}

\noindent
Class records reuse the semiexperimental-fit idea of named parameter classes,
but apply it to diagonal Pulay factors in the frozen GIC basis.  A class has a
name, one factor and one or more case-insensitive label/name patterns separated
by \texttt{|}.  Patterns in a class may deliberately match several GICs.  MATRIX
rejects the class if it matches no GIC or if the matched GICs mix coordinate
families such as stretches, bends, torsions, out-of-plane coordinates, ring
puckering coordinates or special fragment/center coordinates.  This prevents a
chemically motivated scale factor from silently spanning unrelated internal
coordinate types.  Scaling factors must be non-negative.  Ambiguous or
unmatched ordinary selectors remain fatal errors because silent scaling of the
wrong coordinate would invalidate the PED interpretation.

\noindent
Before running GF, the resolved assignments can be inspected without reading an
Hessian and without modifying the xyzin file:

\begin{lstlisting}[style=oracle]
python -m matrix gf \
  --xyzin working/molecule.xyzin \
  --scale-file working/pulay_scale.txt \
  --scale-class "CH_stretches:0.970:R(1,6)|R(2,7)|R(3,8)" \
  --scale-preview
\end{lstlisting}

\noindent
The preview lists every scaling rule, its matched GICs, the detected coordinate
family and the final factor assigned to each GIC after later rules override
earlier ones.

\section{Non-Bonded Hessian Subtraction and Local Models}

Before the Cartesian Hessian is transformed to the GIC basis, GF/PED may
subtract model electrostatic and van der Waals contributions:

\begin{lstlisting}[style=oracle]
python -m matrix gf \
  --fchk working/gauin.fchk \
  --xyzin working/xyzin \
  --subtract-electrostatic \
  --subtract-uff-vdw \
  --nonbonded-14-scale 0.5 \
  --local
\end{lstlisting}

\noindent
The charge vector is read from the canonical \texttt{\#SYNTHONS} section.  If
Gaussian CM5 charges were imported during LINK preprocessing, those are
used.  Otherwise the same synthon charge column is present and comes from the
ORACLE electronegativity model.  GF/PED therefore never reparses Gaussian text
for charges.

The van der Waals correction uses ORACLE topology radii and UFF well-depth
coefficients where available.  Unsupported element pairs fail explicitly
rather than silently using zeros.

Both electrostatic and UFF-vdW corrections use the same topological pair
policy: 1-2 and 1-3 pairs are excluded, 1-4 pairs are included with
\texttt{--nonbonded-14-scale} (default 0.5), and pairs beyond 1-4 are included
with full weight.  Disconnected inter-fragment pairs are treated as beyond
1-4.

The optional \texttt{--local} model is applied after this Cartesian subtraction
and after transformation to internal coordinates.  It keeps diagonal primitive
terms and local bonded couplings while zeroing non-local internal
force-constant terms.

\section{Output Files}

\begin{longtable}{@{}p{0.34\linewidth}p{0.56\linewidth}@{}}
\toprule
Output & Meaning \\
\midrule
\endhead
\texttt{gf\_ped\_report.txt} & Human-readable report for the quick
\texttt{gf} workflow. \\
\texttt{gic\_gf\_ped\_report.txt} & Human-readable report for the frozen-GIC
production workflow. \\
\texttt{gf\_manifest.json} & Run manifest for the quick workflow. \\
\texttt{gic\_gf\_manifest.json} & Run manifest recording FCHK, schema,
optional geometry, scaling input and output files. \\
\texttt{\detokenize{*_frequencies.csv}} & Mode index and frequency in cm\(^{-1}\). \\
\texttt{\detokenize{*_frequency_comparison.csv}} & Optional GF-vs-source
frequency comparison written when the source adapter provides harmonic
frequencies. \\
\texttt{\detokenize{*_gic_labels.csv}} & GIC index, name, irrep and coordinate label. \\
\texttt{\detokenize{*_ped.csv}} & Potential-energy distribution, rows as GICs and columns
as modes. \\
\texttt{\detokenize{*_normal_modes.csv}} & Internal-coordinate normal-mode vectors. \\
\texttt{\detokenize{*_force_constants.csv}} & Internal-coordinate force-constant matrix
after optional scaling. \\
\texttt{\detokenize{*_g_matrix.csv}} & Wilson G matrix in the same GIC order. \\
\texttt{\detokenize{*_g_inverse.csv}} & Global inverse Wilson \(G^{-1}\)
matrix in the same non-redundant GIC order; large-amplitude DVR setup reuses
this table rather than rebuilding reduced moments. \\
\texttt{\detokenize{*_large_amplitude_}}\newline
\texttt{\detokenize{coordinates.csv}} & Large-amplitude GIC
selection with family labels, local one-coordinate GF frequency and
active/excluded status. \\
\texttt{\detokenize{*_large_amplitude_}}\newline
\texttt{\detokenize{blocks.csv}} & Mono- and
poly-dimensional direct \(F/G\) subspace frequencies and coupling
diagnostics, including the combined \(\max(F_{\mathrm{rel}},G_{\mathrm{rel}})\)
criterion and the selected \(G^{-1}_{SS}\) block used by multidimensional DVR
setups. \\
\texttt{\detokenize{*_large_amplitude_}}\newline
\texttt{\detokenize{g_inverse_blocks.csv}} & Long-form
row/column table of the same \(G^{-1}_{SS}\) blocks, convenient for numerical
audit and downstream readers. \\
\texttt{\detokenize{*_large_amplitude_}}\newline
\texttt{\detokenize{mode_ped.csv}} & Total large-amplitude
PED contribution for each GF mode. \\
\texttt{\detokenize{*_large_amplitude_}}\newline
\texttt{\detokenize{dvr_plan.csv}} & DVR-ready
large-amplitude plan with status, \(F/G\) coupling diagnostic, torsional
periodicity, minimum position, \((G^{-1})_{ii}\), and one-term Fourier estimate
when a 1D treatment is justified. \\
\bottomrule
\end{longtable}

The text report begins with the Hessian source, coordinate source, point group,
symmetrization flag and GIC count.  When available, it then prints geometry and
frequency diagnostics before the frequencies, GIC labels and PED matrix.  In
the frozen-GIC workflow, labels and irreps are those stored or evaluated from
the \gicforge{} schema.

\section{GUI Workflow}

The advanced GUI exposes GF/PED through the dedicated \texttt{GF / PED} window.
The window reads a Cartesian Hessian FCHK/FCH file, an optional frozen GIC
definition, an optional B-matrix geometry and optional Pulay scaling records.
It displays the text report and can export the same CSV tables written by the
CLI.

The dashboard also exposes a \texttt{GIC Definition / B Matrix / GF-PED}
workflow.  The intended sequence is:

\begin{enumerate}
  \item define or load a frozen GIC schema;
  \item select the FCHK Hessian;
  \item optionally select a current geometry for B evaluation;
  \item optionally supply Pulay scaling;
  \item run GF/PED and export report/CSV tables.
\end{enumerate}

\section{Validation and Failure Policy}

GF/PED must fail rather than continue when the mathematical model is invalid.
Fatal conditions include:

\begin{itemize}
  \item FCHK missing masses, Cartesian coordinates or Cartesian force constants;
  \item non-symmetric Cartesian Hessian;
  \item schema atom count different from the Hessian input atom count;
  \item B matrix dimension inconsistent with the Hessian;
  \item non-positive-definite Wilson G matrix for the selected coordinates;
  \item non-negligible \(F/G\) couplings between different irreps when
  \texttt{--symmetry-blocks} is requested;
  \item negative Pulay scaling factor;
  \item ambiguous or unmatched Pulay selector.
\end{itemize}

\noindent
A negative vibrational eigenvalue is not a fatal error by itself.  It is
reported as a negative frequency because it may represent a transition-state
direction or an instability in the supplied Hessian/geometry combination.

\section{Relation to GICForge, MORPHEUS and VPT2/VCI}

\gicforge{} owns coordinate construction, pruning and symmetry labels.
GF/PED consumes that frozen coordinate model and performs harmonic vibrational
linear algebra.  MORPHEUS/SEfit consumes the same coordinate model for
semiexperimental structure refinement.  VPT2/VCI is a separate workflow based
on normal-coordinate cubic and quartic force fields and should not be confused
with GIC-based Wilson GF/PED.

This separation is intentional.  A coordinate definition can be generated once,
used to prepare Gaussian input, reused for semiexperimental refinement, and
then reused again for GF/PED analysis of the Cartesian Hessian.  When the same
schema is used throughout, coordinate labels, symmetry species, B matrices and
PED assignments remain auditable across modules.

\end{document}
